% !TEX root = master.tex
% !TEX program = lualatex
\chapter{Review of complex numbers $\mathbb{C}$}\label{chap:}

\section{Basic Definitions}
A complex number is written as
\[
z = x + iy, \qquad x=\Re(z)\in\mathbb{R},\; y=\Im(z)\in\mathbb{R},\; i^2=-1.
\]
\textbf{Conjugate:} $\overline{z}=x-iy$.

\textbf{Modulus:}
\[
|z| = \sqrt{x^2+y^2}.
\]
Geometric interpretation: $|z|$ is the distance from $z$ to $0$, and $|z-z_0|$ is the distance from $z$ to $z_0$.

\section{Euler's Formula}
For $\theta\in\mathbb{R}$,
\[
e^{i\theta}=\cos\theta+i\sin\theta,
\quad\text{thus}\quad |e^{i\theta}|=1.
\]

\section{Polar Coordinates}
If $z\neq 0$, there exist $r=|z|>0$ and an argument $\theta$ such that
\[
z = r e^{i\theta}, \qquad \theta \equiv \Arg(z)\;\;(\mathrm{mod}\;2\pi).
\]
\textbf{Principal argument:} $\Arg(z)\in(-\pi,\pi]$.

\textbf{Products/quotients in polar form:} if $z_1=r_1e^{i\theta_1}$, $z_2=r_2e^{i\theta_2}$,
\[
z_1z_2 = r_1r_2 e^{i(\theta_1+\theta_2)}, \qquad
\frac{z_1}{z_2} = \frac{r_1}{r_2} e^{i(\theta_1-\theta_2)} \quad (z_2\neq 0).
\]
\textbf{Rotation:} multiplying by $e^{i\phi}$ corresponds to a rotation of angle $\phi$.

\section{Algebraic Equations: $k$-th Roots}
To solve $z^k = a$ ($a\neq 0$), write $a=\rho e^{i\alpha}$:
\[
z = \rho^{1/k} e^{i(\alpha+2\pi m)/k}, \qquad m=0,1,\dots,k-1.
\]
Thus there are \textbf{$k$ solutions} (counting multiplicities).
\chapter{Complex functions}\label{chap:}
\section{Complex functions}
A function $f:\Omega\subset\mathbb{C}\to\mathbb{C}$ can be written as
\[
f(z)=f(x+iy)=u(x,y)+iv(x,y),
\]
where $u,v:\Omega\subset\mathbb{R}^2\to\mathbb{R}$.

\subsection{Complex Exponential and Trigonometry}
\[
e^{x+iy}=e^x(\cos y+i\sin y).
\]
Definitions (via the exponential, consistent with Euler’s formula):
\[
\cos z=\frac{e^{iz}+e^{-iz}}{2},\qquad
\sin z=\frac{e^{iz}-e^{-iz}}{2i}.
\]
Hyperbolic functions:
\[
\cosh z=\frac{e^z+e^{-z}}{2},\qquad
\sinh z=\frac{e^z-e^{-z}}{2}.
\]

\section{Continuity}
$f$ is continuous at $z_0$ if
\[
\forall \varepsilon>0,\;\exists \delta>0:\ |z-z_0|<\delta \Rightarrow |f(z)-f(z_0)|<\varepsilon.
\]
Equivalent: $\lim_{z\to z_0} f(z)=f(z_0)$.

\textbf{Practical criterion:} if $f=u+iv$, then $f$ is continuous at $z_0=x_0+iy_0$
\[
\Longleftrightarrow\ u \text{ and } v \text{ are continuous at } (x_0,y_0).
\]

\section{Complex Derivative}
$f$ is \textbf{complex differentiable} at $z_0$ if the limit exists:
\[
f'(z_0)=\lim_{z\to z_0}\frac{f(z)-f(z_0)}{z-z_0}.
\]

\section{Holomorphic and Entire Functions}
Let $\Omega\subset\mathbb{C}$ be \textbf{open}.
\begin{itemize}
\item $f$ is \textbf{holomorphic} on $\Omega$ if it is differentiable at every point of $\Omega$.
\item $f$ is \textbf{entire} if it is holomorphic on $\mathbb{C}$.
\end{itemize}
Recall: $\Omega$ is open if $\forall z_0\in\Omega,\ \exists r>0$ such that the disk
$D_r(z_0)=\{z:|z-z_0|<r\}\subset\Omega$.

\section{Cauchy--Riemann Equations}

\subsection{Theorem (C--R)}
Let $\Omega$ be open and $f=u+iv$ with $u,v\in C^1(\Omega)$.
If $f$ is holomorphic, then
\[
\boxed{\frac{\partial u}{\partial x}=\frac{\partial v}{\partial y}
\qquad\text{and}\qquad
\frac{\partial u}{\partial y}=-\frac{\partial v}{\partial x}}
\quad\text{(Cauchy--Riemann equations).}
\]

\subsection{Useful Formula for $f'$}
When C--R are satisfied,
\[
\boxed{f'(z)=u_x+iv_x=v_y-iu_y.}
\]

\subsection{Quick Examples}
\begin{itemize}
\item $f(z)=z^2=(x+iy)^2=(x^2-y^2)+i(2xy)$:
$u=x^2-y^2,\ v=2xy$.
We have $u_x=2x=v_y$ and $u_y=-2y=-v_x$, hence $f$ is holomorphic and $f'(z)=2z$.
\item Conjugation $f(z)=\overline{z}=x-iy$:
$u=x,\ v=-y$, so $u_x=1$ but $v_y=-1$ $\Rightarrow$ C--R fails:
\textbf{not holomorphic}.
\item $f(z)=e^z=e^x(\cos y+i\sin y)$: C--R holds $\Rightarrow$ holomorphic and $(e^z)'=e^z$.
\item $f(z)=\frac{1}{z}$ is holomorphic on $\mathbb{C}^\ast$ and $f'(z)=-\frac{1}{z^2}$.
\end{itemize}

\subsection{Differentiation Rules (same as real case)}
If $f,g$ are holomorphic on $\Omega$:
\[
(f+g)'=f'+g',\qquad
(fg)'=f'g+fg',\qquad
\left(\frac{f}{g}\right)'=\frac{f'g-fg'}{g^2}\ \ (g\neq 0).
\]
\textbf{Chain rule:} if $g(\Omega)\subset W$ and $f$ is holomorphic on $W$, then
\[
(f\circ g)'(z)=f'(g(z))\,g'(z).
\]
Usual consequences:
\[
(z^k)'=kz^{k-1},\qquad (\cos z)'=-\sin z,\qquad (\sin z)'=\cos z.
\]

\subsection{Geometric Interpretation of C--R (idea)}
View $f:\mathbb{R}^2\to\mathbb{R}^2$ via $(x,y)\mapsto(u(x,y),v(x,y))$.
If $f$ is holomorphic, the Jacobian matrix $Df$ has the form
\[
Df=
\begin{pmatrix}
u_x & u_y\\
v_x & v_y
\end{pmatrix}
=
\begin{pmatrix}
a & -b\\
b & a
\end{pmatrix},
\]
so locally it is a \textbf{scaling} (factor $\sqrt{a^2+b^2}$) followed by a \textbf{rotation}
(no anisotropic stretching): small squares remain squares.

\section{Complex Logarithm}

\subsection{Argument and Discontinuity}
$\Arg(z)$ (principal value) is defined on $\mathbb{C}\setminus\{0\}$ but is \textbf{discontinuous} along the negative real axis
(jump from $\pi$ to $-\pi$).

\subsection{Definition of the Principal Logarithm}
For $z\neq 0$,
\[
\boxed{Log(z)=\log|z|+i\,\Arg(z)}
\]
(principal branch).

\subsection{Properties and Remarks}
\begin{itemize}
\item $Log$ is not defined at $0$ and is not globally continuous (branch cut).
\item On a domain where a branch is chosen (e.g. $\mathbb{C}\setminus\{x\le 0\}$), $Log$ is holomorphic and
\[
(Log z)'=\frac{1}{z}.
\]
\item Usual rules do not hold \emph{strictly} due to multivaluedness/discontinuity:
\[
Log(z_1z_2)\neq Log z_1+Log z_2 \quad \text{in general},
\]
but they hold \textbf{up to a multiple of $2\pi i$}:
\[
Log(z_1z_2)=Log z_1+Log z_2 + 2\pi i k,\qquad k\in\mathbb{Z}.
\]
\end{itemize}
